\documentclass[11pt]{article}
    \title{\textbf{Default Doc}}
    \author{Mint Doan}
    \date{\the\day/\the\month/\the\year}

\input{../pdf.tex}
\input{../webmacros.tex}



\newcommand{\onematrix}{\textbf 1}
\newcommand{\zeromatrix}{\textbf 0}

\usepackage{amsmath}
\begin{document}
 

Before we discuss mathematics properly, since I want this text to be approachable to people with only cursory mathematical experience, I think we should first specify what doing mathematics actually is. There are a few necessary parts of this: why mathematics is possible from its foundations; the structure of mathematics, i.e. what a mathematical discovery would look like if we found one; and finally the language of mathematics, i.e. the way mathematicians articulate our discoveries in order to make them useful. This chapter will be my attempt to discuss these philosophical foundations and motivations, which are usually and unfortunately left implicit by mathematicians; it is generally expected that being able to see these foundations is merely a prerequisite for a practical interest in mathematics. Instead of following this pattern, I will venture to make these prerequisites and expectations explicit, at the risk of putting on the clown makeup for a single chapter.

If you consider yourself sufficiently familiar with the mathematical format, you may wish to skip this chapter, however I think it may hold some value in its articulation of the mathematical pursuit regardless. If you are unfamiliar with the mathematical format, then you should know now that i.e. means "id est", latin for "that is" and generally used to mean "this can also be stated as", and we will be using it a lot.

\section{The Possibility of Mathematics}

First we must discuss "what is mathematics is at its core"? This is a fairly open question, but I think there are two parts of an answer that stand out as insightful: one about the nature of logic that mathematics is built on top of, and one of the mathematical tradition itself. As this project progresses, I hope you will come to see this pattern I will describe which is: that logic, almost alone, has the peculiar power to describe and characterize concepts on an intuitive level, through the network of their consequences and incompatibilities.

The classic example of this is of rain making the ground wet. If it has rained then you know that the ground is wet. If the ground is wet, that does not necessarily imply that has rained. Instead, there could have been some intense humidity, or someone washed their car, or someone simply covered the ground with as much water as they could. If the ground has been wet for a while, then one might deduce that it could not be true that it is a hot dry day, or else the ground would have quickly dried. These statements are incompatible. If I had not told you that the property we were discussing was that "the ground is wet" then from the various things I told you could not be true, or things that would imply this unnamed property, you may even deduce without instruction that we are indeed discussing the proposition "the ground is wet". In the course of 'doing' mathematics, one of the things we do is invent concepts that are unnamed at first, and by their implications, discover what they truly mean.


But where one might say the realm of the mathematician begins is in the full scope of cognitive tools we employ to characterize these propositions. Consider, if we study a two dimensional plane where points have an $x$ and $y$ coordinate so that $(2,3)$ describes a point, I could pose the proposition to you that $y = 5 x + 2$, read as "the $y$ coordinate is equal to five times the $x$ coordinate plus two". Of course this proposition would be wrong in the vast majority of cases, such as $(2,3)$, where $5 x + 2 = 12$ and $y=3$, inequal quantities. However just as we spoke of a space of points in a two dimensional plane, we may perhaps speak of a sub-space where the proposition \emph{is} true, and restrict our attention to this space. If we are to do that, do not need to limit ourselves to discussing this with a mere proposition $y = 5 x + 2$. We can define a set of points where the proposition is true, which we would write as $\{(x,y) \in \reals^2 \hspace{1mm} | \hspace{1mm} y = 5 x + 2\}$, and then we could use the language of set theory to discuss which other sets it intersected with. We could also draw these points on to a two dimensional plane, obtaining a straight line that meets the $x$-axis at $(-0.4,0)$ and passes through at an angle of about 78.69 degrees. On the level of a set, when we show that the intersection of two such sets corresponding to propositions has an empty intersection, we can say that the propositions are mutually exclusive, i.e. if one then not the other. When we draw these graphs, we can check this merely by looking to see if the lines touch.

Mathematics is very much the study of logical abstractions, by restricting them to special cases, seeing which other statements they cannot exist with or must exist with, just as in our example of the ground being wet. However where we break from the study of pure logic is in part that, as mathematicians, we explicitly aim for finding relations and equivalences between different logical groundings, so that we may make deductions using multiple forms of analysis. An equation is a proposition, this is equal to that, but as we saw above it is also a set, vulnerable to set theoretic attack, and also a graph for visual inspection. Moreover, we are actively seeking out new representations to invent and come to understand with new forms of intuition.

One might blame this for the sense that mathematicians quickly devolve into speaking about such foreign concepts, but this is merely the other side of the coin that we are willing to adapt our ways of thinking so drastically to find new ways to understand. In this way and many others, mathematicians sit necessarily between physicists and logicians on the axis of 'logical rigor'. We are distinct from the logicians as we (usually imperceptibly) weaken this rigor to understand things better, and distinct from the physicists as we crave to resolve contradictions in our descriptions that we might finally 'know what we are talking about'.

\section{The Structure and Utility of a (Mathematical) Theory}

To describe what one hopes to achieve from mathematics, it's perhaps easier to first describe what is achieved from a theory in abstract. When one selects an object of study, say for example, mammals, one may first put significant effort into defining the object, and then giving structure to subcategories in order to discover meaning from those substructures.

I am no biologist but naively and for the sake of argument, one may say that it is clear a mouse is a mammal, and clear that a lizard is not. One may then decide that a bird is not a mammal, and thus that mammals give live birth and do not have beaks. However many echidnas have beaks, and platypus have beaks as well as laying eggs. Despite this exception, they both still have only one jaw bone, do not have feathers, and lactate to feed their newborns, so we may broaden or adjust our conditions slightly and still call them mammals. With the boundaries around our study fixed, we may then notice that bears and dogs share a resemblance, and that dogs and cats share a similar body plan, constructing the order Carnivora and the two suborders Feliformia and Caniformia, containing cats and dogs together with bears respectively. Such a structure is necessary for one to argue for Speciation, that just as we have artificially bred kinds of dogs, that perhaps nature has over a long time bred some four legged animal into what became a proto-dog and proto-cat, and the former was later separated into dogs and bears. Now if you were responsible for medically treating one of these animals, you would have the information that the biology of a dog is more similar to a bear than to a cat. And you can conceive of a broader family-tree of life for which you may extend this utility or draw on further examples to better inform it.

That contrived example now concluded, let me describe how this works for something such as Group theory.

For the layman with great interest in mathematics, a naive attempt to learn about it may yield what a group is axiomatically, a set closed under a multiplication operation which is associative such that every element has a multiplicative inverse and an identity element, or an explanation grounded in intuition such as the ways in which a shape can be rotated or moved so that it looks the same, or algebraic examples such as the set of matrices with non-zero determinants. To use our previous example, this would be like noticing the existence of the category of mammals and stopping there. It is not clear what deeper understanding we have achieved by constructing this, we have not discovered any particular similarity between groups nor do we have a notion of a family tree with which to discover other valuable patterns.

So let me describe how that works.

Once we have two groups, $G$ and $H$, we may construct new groups, such as the product group $G \times H$, whose members are pairs $(g,h)$ for each $g$ in $G$ and $h$ in $H$. If we have a multiplication operation between groups, can we reverse this? Is there a notion of group division?

In fact there is. First we find a subgroup $H$ to divide, that is, a subset of $G$ inheriting the multiplication operator, which cannot be escaped into $G$ by multiplying its elements and also contains all of its own inverses. A group is not always commutative, that is, we cannot say for certain that two elements $g$ and $h$ in $G$ satisfy $g h = h g$, i.e. multiplication cannot be reordered arbitrarily. However there may be subgroups whose elements do reorder with the rest of the group, i.e. a subgroup $Z$ in $G$ for which all elements $z$ satisfy $g z = z g$ for $g$ in $G$ even if this is not satisfied by other elements $h$ in $G$ (we would call $Z$ the 'center' of $G$). More generally, a subgroup $N$ may not necessarily have elements that each reorder with each element of $G$, but may reorder as a whole subgroup, so for example, for each $g$ in $G$, we have $g h = k g$ where $h eq.not k$ most of the time but $h$ and $k$ are both in $N$. So long as this is true, we write $g N = N g$ to mean 'some element of $N$' in place of the element itself, i.e. for any $h$ in $N$, there will be some $k$ in $N$ such that $g h = k g$, and when this is true, we call $N$ a 'normal' subgroup. If we have $g N = N g$, then elements of $N$ may be in a sense reordered with the rest of the elements in $G$, and so we may define a new group $G/N$ of elements $g N$ for each $g$ in $G$. This 'quotient group' factors out the dynamics of the normal group, since we have $(g N) (h N) = g (N h) N = g (h N) N = (g h) N$ for elements $g$ and $h$ in $G$, and should $h$ also be a member of the subgroup $N$, then we have $g (h N) = g N$, meaning that any elements of $N$ are factored out.

It is not important to understand this example in detail, in fact later we will discuss it properly, however the example does illustrate the point. We take a definition of a group and construct a group product, and a subgroup, then the conditions necessary for a factor group and thus a quotient group. If we restrict our attention to the finite groups, then we may ask whether or not this notion of divisibility admits a version of 'prime numbers' for groups, and indeed this is true. In fact, mathematicians have completed a de facto periodic table of finite groups, the classification of finite simple groups, and proved from underlying axioms that all possible finite groups are accounted for, and that it is impossible for a group to exist which does not classify into our periodic table.

From here, the intuition that groups describe "symmetries" of objects now has strong rules: the symmetries of objects decompose into now well known groups and we may say with certainty what the atomic components of these symmetries are. Our theory has yielded an unfalsifiable taxonomy and practical implications for all systems with discrete symmetries.

Many fields in mathematics follow a similar pattern as I laid out in my example about mammals which we then applied to group theory. The protocol, if you can even call it that, is thus: strictly define a core object, preferably abstractly so that later on you can apply the structure to anything that fits the bill (e.g. symmetries described by group theory); create sub-labels (e.g. subgroups, some subgroups are 'normal' or 'central') and characterize their properties to build an understanding of what these labels are are or what they mean, possibly relating the structure to something more familiar (e.g. product groups and quotient groups as analogous to multiplication and division of numbers); and finally, substantiate the utility of this theory if you have not already, which we did by saying that 'groups describe symmetries', which is formalized in the study of 'group actions'.

With this pattern in mind, we are able to collect  'theories' or academic sub-disciplines centered around a school of thought, each with a certain descriptive capability. As I said in step one of the protocol, this description is abstract, so one may levy the accusation that group theory describes only the theory of groups, except that in step three we substantiate that groups are everywhere. With a large collection of theories each studying some construction that 'appears everywhere', we assemble an increasingly comprehensive framework that is able to provide already well studied models for new phenomena that we discover. For instance, once the differential equation is well studied, the experimental evidence that not a magnetic field but a \emph{change} in magnetic fields induces a current, allows the articulation of Faraday's law. Or we could apply our models retroactively to draw new conclusions: man has always known that fluid flows and perhaps at times even understood many things about its flow, but only with an understanding of vector calculus and continuum mechanics did it become reasonable to describe Navier-Stokes equations, and only much later did we discover a framework by which to use this model computationally (e.g. finite volume methods) when we could not solve the equations analytically.

\section{What is Doing Mathematics vs Using Mathematics}


 
\end{document}